\documentclass[11pt]{article}
\usepackage[letterpaper,margin=0.88in]{geometry}
\usepackage[T1]{fontenc}
\usepackage{lmodern}
\usepackage{microtype}
\usepackage{amsmath,amssymb,amsthm,booktabs,array}
\usepackage{xcolor,fancyhdr}
\usepackage[colorlinks=true,linkcolor=blue!45!black,citecolor=blue!45!black,urlcolor=blue!45!black]{hyperref}
\definecolor{ink}{HTML}{182838}
\definecolor{rulegray}{HTML}{CCD4DD}
\color{ink}
\hypersetup{pdftitle={An explicit depth bound for sorting networks},pdfsubject={Finite parameters for Chvatal's AKS construction}}
\input{glyphtounicode}
\pdfgentounicode=1
\pagestyle{fancy}
\fancyhf{}
\fancyhead[L]{\small\scshape Sorting networks}
\fancyhead[R]{\small\thepage}
\renewcommand{\headrulewidth}{0.4pt}
\setlength{\headheight}{14pt}
\setlength{\parindent}{0pt}
\setlength{\parskip}{5pt plus 1pt}
\setlength{\abovedisplayskip}{8pt plus 2pt minus 2pt}
\setlength{\belowdisplayskip}{8pt plus 2pt minus 2pt}
\setlength{\abovedisplayshortskip}{5pt plus 2pt}
\setlength{\belowdisplayshortskip}{5pt plus 2pt}
\raggedbottom
\newcommand{\lb}{\log_2}
\newcommand{\Eb}{\varepsilon_B}
\newcommand{\Ef}{\varepsilon_F}
\newcommand{\df}{\delta_F}
\newcommand{\depth}{\mathcal{L}}
\newcommand{\bc}[1]{B(#1)}
\begin{document}
\thispagestyle{plain}
{\LARGE\bfseries An explicit depth bound for sorting networks\par}
\vspace{4pt}
{\large A finite parameter choice for Chv\'atal's AKS construction\par}
\vspace{7pt}
\noindent\textbf{Theorem.}
Let \(D(N)\) be the minimum depth of a comparator sorting network on \(N\) wires.
For every integer \(N\ge 2^{79}\),
\[
  \boxed{D(N)\le \frac{5236}{3}\lb N-40651.}
\]
The coefficient \(5236/3\approx 1745.33\) is about \(4.63\%\) smaller than the coefficient \(1830\) in Chv\'atal's Theorem~1.1~\cite{chvatal}. The proof uses the separator construction and outsider invariant from that report, with all parameters fixed below.

\section{Separator sizes}

First take \(N=k^d\) with \(d\ge14\). Throughout, let
\[
 k=64,\qquad A=2048,\qquad \nu=2^{-5},\qquad M=2^{59}.
\]
Write \(B(p)=p(p+1)/2\) for the depth of Batcher's odd--even mergesort on \(2^p\) wires~\cite{batcher}. The same depth is an upper bound for smaller inputs by padding with largest keys and deleting their fixed paths. In Chv\'atal's construction, a non-root separator on \(a\le M\) wires can be implemented by a full sorter. It has depth at most \(B(59)=1770\). At a bottom level that sends all of its wires to the parent, the construction needs only routing.

For the remaining non-root separators, write the fringe block sizes and the \(k\) middle block sizes as \(2^s f_0\) and \(2^s b_0\), respectively, in the notation of \cite[\S5]{chvatal}. Thus \(a=2^s m_0\), where \(m_0=2f_0+kb_0\). Choose the largest integer \(r\) with \(2^r m_0\le M/16\), and set
\[
 m=2^r m_0,\qquad n=2^{s-r},\qquad f=2^r f_0,\qquad b=2^r b_0.
\]
Here \(M\ge32A^2k^2=2^{39}\), and the cases below all satisfy \(m_0<2A^2k^2\). Since \(a>M\), it follows that \(s>r\) and \(n=a/m>16\).

\begin{center}
\small
\setlength{\tabcolsep}{9pt}
\begin{tabular}{@{}lrrrr@{}}
\toprule
Separator position & \(m_0\) & \(r\) & \(m\) & \(f\) \\
\midrule
Top, advancing & \(2^6\) & \(49\) & \(2^{55}\) & \(0\) \\
Top, retreating & \(2^{28}\) & \(27\) & \(2^{55}\) & \(2^{32}\) \\
Interior & \(2^{35}-2\) & \(20\) & \(2^{55}-2^{21}\) & \(4095\cdot2^{20}\) \\
Bottom, advancing & \(2^{35-h}-2\) & \(20+h\) & \(2^{55}-2^{21+h}\) & \(4095\cdot2^{20+h}\) \\
\bottomrule
\end{tabular}
\end{center}

At an advancing bottom level \(i\) at time \(t\), define \(2^h=c(i,t)/(Nk^{-i})\), where \(c(i,t)=(N/(A\nu k))A^i\nu^t\). The bottom-level capacity bound gives \(h\ge0\). Because the bottom advances, \(c(i,t)=c(i+1,t+1)/(A\nu)<(Ak/\nu)Nk^{-i}\), so \(h\le21\). Substituting this ratio into the block sizes of \cite[\S5]{chvatal} gives the last row. The top-retreating row also applies at \(t=1\), by the explicit transfer from level~1 in \cite[\S3]{chvatal}. Consequently every non-root scramble has \(m\ge m_{\min}=2^{55}-2^{42}\). Every nonempty fringe is even and has \(f\ge f_{\min}=4095\cdot2^{20}\).

\newpage
\section{Separator quality and the outsider invariant}

Use the exact rational parameters
\[
 \mu=2^{-28},\qquad \df=\frac{128}{4095},\qquad
 \Ef=\frac{1}{20{,}000{,}000},\qquad
 \Eb=\frac{467}{10^{10}},\qquad \delta=\frac{8}{10^{10}}.
\]
Theorem~5.1 of \cite{chvatal} applies to each of the non-root scramble shapes with nonempty fringe. The bulk-error requirement decreases with \(m\), and
\[
 \frac{2(1+\ln m)}{m}\le
 \frac{2(1+\ln m_{\min})}{m_{\min}}
 < \frac{2(1+55\cdot 0.694)}{m_{\min}} < \Eb^2.
\]
For the fringe, \(\df<1/25\) and \(1/\Ef\) is an integer. Put \(L=\ln(0.12/(e\df))\). We have \(L>1/3\) and \(\ln(3e^5 f_{\min})<35\). The function
\[
 f\longmapsto \frac{2}{f-2}\left(1+\frac{\ln(3e^5f)}{L}\right)
\]
is decreasing for \(f\ge f_{\min}\), as follows by differentiation. Therefore, for each nonempty fringe,
\[
 \frac{2}{f-2}\left(1+\frac{\ln(3e^5f)}{L}\right)
 < \frac{212}{f_{\min}-2}<\Ef,
 \qquad \frac{4e}{f}<\frac{12}{f_{\min}}<\Ef.
\]
The bound \(\ln 2<0.694\) follows by summing four terms of the series \(2\sum_{j\ge0}3^{-(2j+1)}/(2j+1)\) and bounding the tail by a geometric series. The Taylor series for \(e^x\) gives \(e^{4/3}<19/5<0.12/\df\) by summing through degree~10 and bounding the tail by its first term times \(9/8\). Summing through degree~45 at \(x=35\) gives \(e^{35}>300{,}000 f_{\min}>3 e^5 f_{\min}\), using \(e<3\).

When \(f=0\), the fringe condition is vacuous and Lemma~6.1 of \cite{chvatal} gives Property~B, the required bound on bulk errors. In particular, use a separate root scramble with \(m^*=2^{75}\) and \(n^*=2^{6d-75}\ge16\). Its \(k\) equal output blocks each consist of \(m^*/k\) rows. Take \(\varepsilon^*=\mu/k=2^{-34}\). Its bulk-error requirement holds since
\[
 \frac{2(1+75\ln 2)}{2^{75}}<
 \frac{2(1+75\cdot0.694)}{2^{75}}<(2^{-34})^2.
\]

The remaining conditions are those of \cite[\S4]{chvatal}. Condition (4.1) holds with equality by the choice of \(\varepsilon^*\). Conditions (4.3) and (4.4) hold with equality since \(\mu=\nu/(Ak^2)=\df(A\nu k-1)/(2A^2k^2)\). The denominator in (4.2) is positive, since \(\delta kA<1\), and rational substitution bounds the normalized left-hand side of (4.2) and the left-hand side of (4.5), respectively
\begin{align*}
 \frac{1}{A\nu}\left[1+\frac{(k-1)\delta k A^2}{1-\delta^2 k^2 A^2}
   +\frac{A\nu k-2A\nu+1}{2\mu A^2 k^2}
   +\frac{\Eb}{\mu}\right]+\frac{\delta Ak}{\nu}
   &< \frac{911}{1000},\\
 \frac{\Ef}{A\nu}+\frac{\delta^2Ak}{\nu}
   &=7.8393435456\times10^{-10}<\delta.
\end{align*}
At the stopping time, Chv\'atal's capacity bound is \(c(\alpha(t_f),t_f)<2A^2k^2=2^{35}\), and \(\mu\delta\,2^{35}=1.024\times10^{-7}<1\). Lemma~4.5 thus eliminates all outsiders of order~1. Every key at the stopping level is addressed to a leaf below its node's parent. Sorting all the wires grouped under each such parent completes the network.

\newpage
\section{Stopping time and depth}

For these parameters the unrounded top and bottom levels of \cite[\S3]{chvatal} are
\[
 \alpha^*(t)=\frac{5t-6d+25}{11},\qquad
 \omega^*(t)=\frac{5t+12}{17}.
\]
Each is rounded up to an integer of parity \(t\), with the top also constrained to be nonnegative. The stopping time \(t_f\) is the first \(t\ge2\) at which the rounded levels coincide. At a meeting of level \(i\), set \(q=d-i\) and \(u=5t-17i\). The bottom is at least~2 for \(t\ge2\), so the nonnegative constraint does not alter the meeting condition. In terms of \(q\) and \(u\), the meeting conditions are
\[
 \max\{-46,6q-47\}<u\le\min\{-12,6q-25\},
 \qquad
 u\equiv -12(d-q)\pmod{10}.
\]
The congruence enforces that \(t=(17i+u)/5\) is an integer with the same parity as \(i\). The interval requires \(-3\le q\le5\). If \(q=3\) then \(u\ge-28\), and if \(q\le2\) then \(u\ge-45\). In either case \(5t-17d=u-17q\ge-79\). For \(q=4\) the interval for \(u\) is \([-22,-12]\), and for \(q=5\) it is \([-16,-12]\). Taking the least admissible \(u\) in each interval for each residue class of \(d\), then minimizing \(u-17q\), gives the table below. Every class has a meeting with \(u-17q<-79\), so these are the first meetings. Each row extends from its representative \(d\) to \(d+5j\) by increasing \(t_f\) by \(17j\) and keeping \(q\) fixed.

\begin{center}
\small
\setlength{\tabcolsep}{9pt}
\begin{tabular}{@{}rrrrr@{}}
\toprule
Representative \(d\) & \(t_f\) & \(q=d-i\) & \(5t_f-17d\) & \(\depth(d)-10472d\) \\
\midrule
14 & 30 & 4 & \(-88\) & \(-51123\) \\
15 & 33 & 4 & \(-90\) & \(-52355\) \\
16 & 35 & 5 & \(-97\) & \(-56466\) \\
17 & 38 & 5 & \(-99\) & \(-57698\) \\
18 & 41 & 5 & \(-101\) & \(-58930\) \\
\bottomrule
\end{tabular}
\end{center}

The root uses two rounds of \(2^{75}\)-sorters. Each subsequent scramble uses two rounds on at most \(2^{55}\) wires per column, so \(2B(55)=3080\) layers suffice per step, including steps that use full sorters. After step \(t_f\), a parent on level \(i-1\) has \(k^{q+1}\) wires below it. All parents can sort in parallel. The resulting depth is therefore at most
\[
 \depth(d)=2B(75)+(t_f-1)\,2B(55)+B\bigl(6(q+1)\bigr).
\]
The maximum in the last column is \(-51123\). Since \(d\mapsto d+5\) adds exactly \(17\cdot3080=5\cdot10472\) to \(\depth(d)\), the table proves \(D(64^d)\le10472d-51123\) for all \(d\ge14\). For arbitrary \(N\ge2^{79}\), pad to \(64^d\) with \(d=\lceil\lb N/6\rceil\ge14\). Then \(10472d\le(5236/3)\lb N+10472\), which gives the stated bound. \hfill\(\square\)

\begin{thebibliography}{9}
\setlength{\itemsep}{1pt}
\small
\bibitem{chvatal} V. Chv\'atal. \emph{Lecture Notes on the New AKS Sorting Network}. Rutgers University, DCS-TR-294, 1992. \href{https://doi.org/10.7282/T379487B}{doi:10.7282/T379487B}.
\bibitem{batcher} K. E. Batcher. Sorting networks and their applications. \emph{AFIPS Spring Joint Computer Conference}, pp.~307--314, 1968. \href{https://doi.org/10.1145/1468075.1468121}{doi:10.1145/1468075.1468121}.
\end{thebibliography}
\end{document}
